\subsection{Code Structure}
The controller is written in C++. It does not call board functions, so the same
code can be tested on a computer. The sketch only reads inputs, calls the
controller and writes outputs. The browser model follows the same state rules.

\begin{tabularx}{\linewidth}{@{}p{5.6cm}Y@{}}
\toprule
File & Purpose\\
\midrule
\code{controller.h} & States, credit, timer, STOP count, outputs and debounce.\\
\code{washing\_controller.ino} & Pin setup, physical input reading and Serial test commands.\\
\code{simulator/index.html} & Offline page with coins, buttons, LEDs, timer and event log.\\
\code{simulator/controller.js} & Browser version of the state controller.\\
\code{tests/test\_controller.py} & Compiles the C++ tests and checks expected results.\\
\code{tests/test\_firmware.cpp} & Runs the actual sketch with mocked board inputs and outputs.\\
\bottomrule
\end{tabularx}

\subsection{Important Code}
The deadline is checked before control buttons. Both active states use the same
start time. These lines are taken from the controller:
\begin{lstlisting}
if (active() && uint32_t(now - started_) >= CycleMs) {
    finish(now);
    return;
}
if (in.stop) {
    ++stop_count_;
    if (stop_count_ >= 2) finish(now);
    return;
}
if (in.pause && state_ == State::Running) {
    enter(State::Paused, now);
} else if (in.run && state_ == State::Paused) {
    // Keep started_: a pause never gives extra time.
    enter(State::Running, now);
}
\end{lstlisting}
The sketch reads \code{millis()} and uses unsigned subtraction for elapsed time.
It updates the timer and LEDs without \code{delay()}, following the idea of
Arduino's Blink Without Delay example~\cite{blink}. The tests also cover clock
rollover and button bounce.

\subsection{Running the Demonstration}
Open \code{simulator/index.html} in a browser. Insert coins, press RUN, advance
eight minutes and press PAUSE. Advance another eight minutes and press RUN.
The display should show 14:00 remaining. Advance to minute 30 to verify timeout.
Test two STOP presses in a new cycle, then inject a fault.

The default clock is manual. Real time runs at 1x; demo time runs at 60x. Both
use the same 30 minute virtual cycle. The browser can export the event log.

For the board sketch, open \code{washing\_controller.ino} with its adjacent
\code{controller.h}. Serial commands are \code{1}, \code{2}, \code{5} for coins
and \code{r}, \code{p}, \code{s} for controls at 115200 baud. Verify and upload
to the chosen board before using it in a physical demonstration.
